\documentclass[tikz,border=4pt]{standalone}
\usepackage{ctex}
\usepackage{pgfplots}
\usepgfplotslibrary{groupplots}
\pgfplotsset{compat=1.18}
\begin{document}
\begin{tikzpicture}[
    every axis/.style={width=5cm, height=4cm, axis lines=middle, samples=200, grid=both, grid style={gray!20}, title style={yshift=-1ex}},
]
\begin{groupplot}[
    group style={group size=3 by 1, horizontal sep=1.5cm},
]
% (1) 可去间断 f(x) = (x^2-1)/(x-1), x=1 处 = 2 但定义 f(1)=0
\nextgroupplot[xmin=-1, xmax=3, ymin=-1, ymax=3, title={可去间断}]
\addplot[blue, thick, domain=-1:0.95] {x+1};
\addplot[blue, thick, domain=1.05:3] {x+1};
\fill[white] (axis cs:1, 2) circle (2.5pt);
\draw[blue] (axis cs:1, 2) circle (2.5pt);
\fill[blue] (axis cs:1, 0) circle (2pt);
\node[anchor=west, font=\tiny] at (axis cs:1.1, 0.3) {$f(1)=0$};

% (2) 跳跃间断
\nextgroupplot[xmin=-1, xmax=3, ymin=-1, ymax=3, title={跳跃间断}]
\addplot[blue, thick, domain=-1:1] {x};
\addplot[blue, thick, domain=1:3] {x-1};
\fill[blue] (axis cs:1, 1) circle (2pt);
\fill[white] (axis cs:1, 0) circle (2.5pt);
\draw[blue] (axis cs:1, 0) circle (2.5pt);

% (3) 第二类（无穷型）
\nextgroupplot[xmin=-1, xmax=3, ymin=-3, ymax=3, title={第二类（无穷型）}]
\addplot[blue, thick, domain=-1:0.95, restrict y to domain=-3:3] {1/(x-1)};
\addplot[blue, thick, domain=1.05:3, restrict y to domain=-3:3] {1/(x-1)};
\draw[red, dashed] (axis cs:1, -3) -- (axis cs:1, 3);
\end{groupplot}
\end{tikzpicture}
\end{document}
